\newglossaryentry{bitexte}{
    name=bitexte,
    description={
        «~Dans le domaine de la traduction, un bitexte est un document fusionné composé de la version dans une
        langue source et une langue cible d'un texte donné.
        Les bitextes sont générés par un logiciel appelé un outil d'alignement, ou outil bitexte, qui aligne
        automatiquement les versions originales et traduites d'un même texte. L'outil fait correspondre ces
        deux textes phrase par phrase.~»
        (\href{https://en.wikipedia.org/wiki/Parallel_text}{Wikipedia/Parallel\_text dans la partie "Bitext"})\footnote{«~In the field of translation studies a bitext is a merged document
        composed of both source- and target-language versions of a given text. Bitexts are generated by a piece of software
        called an alignment tool, or a bitext tool, which automatically aligns the original and translated versions of the same
        text. The tool generally matches these two texts sentence by sentence.~» notre tradcution}
    }
}

\newglossaryentry{BLEU}{
    name=BLEU,
    description={
        BLEU (\textit{BiLingual Evaluation Understudy}) est un algorithme d’évaluation de la qualité d'un
        texte traduit automatiquement d’une langue naturelle à une autre. La qualité est considérée
        comme la correspondance entre la production d’une machine et celle d’un humain. Le score est compris
        entre 0 et 1, 0 indiquant une traduction éloignée d'une traduction humaine, et 1 indiquant une
        traduction automatique très similaire. (\href{https://docs.cloud.google.com/translate/docs/bleu-scores?hl=fr}{Évaluer des modèles par Google})
    }
}

\newglossaryentry{CAHAI}{
    name=CAHAI,
    description={
        Le CAHAI ou Comité \textit{ad hoc} sur l’intelligence artificielle est un comité ayant été en
        mandat de 2019 à 2021, celui-ci examinait «~la faisabilité et les éléments potentiels, sur la base de
        larges consultations multipartites, d'un cadre juridique pour le développement, la conception et
        l'application de l'intelligence artificielle, fondé sur les normes du Conseil de l'Europe en matière
        de droits de l'Homme, de démocratie et d'État de droit~». Il a depuis été remplacé par le
        Comité sur l'IA (CAI). (\href{https://www.coe.int/fr/web/artificial-intelligence/cahai}{site du Conseil de l'Europe (COE) sur le CAHAI})
    }
}

\newglossaryentry{CLIR}{
    name=CLIR,
    description={
        La recherche d'informations interlangue ou \textit{Cross-language information retrieval} (CLIR) est une tâche
        de recherche dans laquelle les requêtes de recherche et les documents candidats sont rédigés dans des
        langues différentes. La CLIR peut être très utile dans certains scénarios. Par exemple,
        un journaliste peut vouloir rechercher des informations en langue étrangère pour obtenir différentes
        perspectives pour son article ; un inventeur peut explorer les brevets d'un autre pays pour comprendre
        l'état de l'art. (\href{https://en.wikipedia.org/wiki/Cross-language_information_retrieval}{Wikipedia/Cross-language\_information\_retrieval})
    }
}

\newglossaryentry{CNN}{
    name=CNN,
    description={
        «~En apprentissage automatique, un réseau de neurones convolutifs ou réseau de neurones
        à convolution (en anglais CNN ou ConvNet pour \textit{convolutional neural network}) est un type de réseau
        de neurones artificiels acycliques (\textit{feed-forward}) inventé en 1980 par Kunihiko Fukushima,
        dans lequel le motif de connexion entre les neurones est inspiré par le cortex visuel des animaux.~»
        (\href{https://fr.wikipedia.org/wiki/Réseau_neuronal_convolutif}{Wikipedia/Réseau\_neuronal\_convolutif})
    }
}

\newglossaryentry{fonction softmax}{
    name=softmax,
    description={
        La fonction softmax est une généralisation de la fonction sigmoïde aux cas multi-classes.
        Elle prend en entrée un vecteur de valeurs réelles et le transforme en une distribution de
        probabilités, où chaque valeur de sortie est comprise entre 0 et 1 et où la somme de
        l'ensemble des valeurs vaut exactement 1. Dans le cadre des réseaux de neurones, elle est
        typiquement utilisée en couche de sortie pour des tâches de classification
        \citep{soulie_1990}.
    }
}

\newglossaryentry{régression logistique}{
    name=régression logistique,
    description={
        «~Il s'agit d'une technique de modélisation qui [...] vise à prédire et expliquer les valeurs
        d'une variable catégorielle binaire $Y$ à partir d'une collection de variables
        $X$ continues ou binaires~» \citep{rakotomalala_2015}.
    }
}

\newglossaryentry{MNRL}{
    name=MNRL,
    description={
        Le principe MNRL (\textit{Multiple Negative Ranking Loss}) est de rendre claire la distinction entre
        des données pertinentes (positives) et non pertinentes (négatives) \citep{henderson_2017}.
    }
}

\newglossaryentry{n-gramme}{
    name=n-gramme,
    description={
        Un n-gramme est une suite de n éléments consécutifs extraite d’un texte ou d’une séquence.
        Les éléments peuvent être des mots, des caractères, voire des tokens selon l’application.
    }
}

\newglossaryentry{plongements lexicaux}{
    name=plongements lexicaux,
    description={
        «~Les plongements lexicaux ou plongements sémantiques ({« \textit{word embedding} »} en anglais) est une méthode
        d'apprentissage d'une représentation de mots sous forme de vecteurs, utilisée notamment en
        traitement automatique des langues.~»
        (\href{https://fr.wikipedia.org/wiki/Plongement_lexical}{Wikipedia/Plongement\_lexical})
    }
}

\newglossaryentry{plongement de phrases}{
    name=plongement de phrases,
    description={
        «~En traitement automatique des langues, le plongement de phrases est une représentation
        d'une phrase sous forme de vecteurs de nombres qui code des informations sémantiques significatives.~»
        (\href{https://en.wikipedia.org/wiki/Sentence_embedding}{Wikipedia/Sentence\_embedding}) \footnote{
        «~In natural language processing, a sentence embedding is a representation
        of a sentence as a vector of numbers which encodes meaningful semantic
        information.~» notre traduction}
    }
}

\newglossaryentry{réseau siamois}{
    name=réseau siamois,
    description={
        Le réseau siamois se compose de deux \gls{CNN} et utilise une perte contrastive.
        Le réseau Triplet se compose de trois CNN et utilise une perte triplet \citep{zhang_2020}.
    }
}

\newglossaryentry{similarité cosinus}{
    name=similarité cosinus,
    description={
        Un calcul de similarité entre deux vecteurs à \textit{n} dimensions en déterminant le cosinus de
        leur angle. «~La valeur d'un cosinus [...], est comprise dans l'intervalle [-1,~1]. La valeur de -1 indique
        des vecteurs opposés, la valeur de 0 des vecteurs indépendants (orthogonaux) et la valeur de 1 des
        vecteurs colinéaires de coefficient positif. Les valeurs intermédiaires permettent d'évaluer le degré
        de similarité.~» (\href{https://fr.wikipedia.org/wiki/Similarité_cosinus}{Wikipedia/Similarité\_cosinus}).
    }
}

\newglossaryentry{similarité sémantique}{
    name=similarité sémantique,
    description={
        «~La similarité sémantique est une notion définie entre deux concepts, soit au sein d'une
        même hiérarchie conceptuelle, soit - dans le cas d'alignement d'ontologies - entre deux concepts 
        appartenant respectivement à deux hiérarchies conceptuelles distinctes. La similarité sémantique
        indique que ces deux concepts possèdent un grand nombre d'éléments en commun (propriétés, termes,
        instances).~» (\href{https://fr.wikipedia.org/wiki/Similarité_sémantique}{Wikipedia/Similarité\_sémantique})
    }
}

\newglossaryentry{transformers}{
    name=transformers,
    description={
        Un transformer est une architecture d'algorithme neuronal.
        Les transformers sont conçus pour gérer des données séquentielles, notamment du texte,
        pour des tâches telles que la traduction et la génération de texte.
        Le terme est aussi utilisé en français \citep{yvon_2022}.
        (\href{https://fr.wikipedia.org/wiki/Transformeur}{Wikipedia/Transformeur})
    }
}

\newglossaryentry{courbes AUC-ROC}{
    name=courbes AUC-ROC,
    description={
        «~La courbe ROC (Receiver Operating Characteristic) est tracée en calculant le taux de vrais positifs (TVP)
        et le taux de faux positifs (TFP) à chaque seuil possible (en pratique,
        à des intervalles sélectionnés), puis en représentant le TVP par rapport
        au TFP. Un modèle parfait, qui, à un certain seuil, présente un TPR de
        1,0 et un FPR de 0,0, peut être représenté par un point en (0, 1) si
        tous les autres seuils sont ignorés.~»\\
        «~L'aire sous la courbe ROC (AUC ou \textit{area under curve}) représente la probabilité que le modèle, s'il reçoit un exemple positif
        et un exemple négatif choisis aléatoirement, classe l'exemple positif plus haut que l'exemple négatif.~»
        (\href{https://developers.google.com/machine-learning/crash-course/classification/roc-and-auc?hl=fr}{Google for developers})
    }
}

\makenoidxglossaries